\justifying
\NSFCBodyText

\subsubsection{已具备的实验条件}
申请人依托新疆大学开展研究。已确认可稳定使用的显卡包括 NVIDIA A40 8块和 NVIDIA GeForce RTX 3090 Ti 6块。现有软件工具包括 Linux、Python、深度学习框架、Git 和 LaTeX 等。这些条件能够支持数据质量检查、多尺度建模、概率预测训练和滚动评价。研究拟采用有限规模模型与统一计算预算，优先保证基线、消融和独立测试的完整性。

已有光伏功率预测和长时序预测论文提供了方法研究基础。申请人将负责总体方案和实验结果核查，具体参与者及任务分配见\textbf{【待补 F-XJ-03：本项目人员及投入安排】}。学院公共算力可作为拟申请的补充条件，现阶段不计入固定资源配额。

当前材料没有给出已获授权的新疆光伏场站实测数据、发布时刻可追溯的业务 NWP 档案或本项目专属数据库。公开光伏与气象数据拟用于前期方法验证，正式使用前需核对数据许可、时间覆盖与变量定义，并保留来源和版本。公开数据上的结果与新疆实地检验分别记录，避免扩大结论适用范围。

\subsubsection{尚缺少的实验条件及拟解决的途径}
\textbf{新疆光伏场站实测数据。}拟获取能够关联功率、时间戳、装机容量与运行状态的历史记录，并核对缺测、停机和限电等标记。数据提供单位、授权范围、采样频率及覆盖时段为\textbf{【待补 F-XJ-01：新疆场站数据与使用授权】}。拟由申请人在项目启动阶段对接数据权利人，先完成数据清单与质量评估，再确定可开展的预测时间尺度和天气分层分析范围。尚未取得的数据不计入已有实验条件。

\textbf{可追溯的气象预报输入。}业务 NWP 的来源、发布时刻、预报步长和档案范围为\textbf{【待补 F-XJ-02：业务 NWP 档案与获取条件】}。拟以实际预测时刻能够获得的信息组织回溯实验。再分析资料可支持背景描述和机制探索，但不作为已发布的业务预报替代物；若无法取得相应档案，则明确将该部分结果限定为离线方法研究，不据此声称达到业务运行水平。

\textbf{合作验证与人员安排。}拟通过气象业务单位或场站获取需求反馈和案例核验，具体合作方、承担事项与书面落实情况为\textbf{【待补 F-XJ-04：合作单位及协议落实情况】}。人员分工按实际加入本项目的成员确定，不沿用其他申请的团队名单。建议经费3万元主要用于与数据整理、实验验证及成果组织相关的必要支出，细项按实际适用的基金管理要求另行核定；现有计算资源优先用于模型试验。

\subsubsection{利用研究基地的计划与落实情况}
本项目拟以申请人已确认可用的校内研究条件为基础开展计算实验，当前未提供国家实验室、全国重点实验室、部门重点实验室或相关业务基地的正式接入与合作证明。若后续需要依托气象业务研究基地获取数据或开展评价，将按实际协议明确数据用途、验证任务和成果使用范围，并据实更新落实状态。基地支持条件与前述合作缺口一并核实，不将意向接洽表述为已落实支撑。
